\chapter{Composite Bloch Frames, Gauge Alignment, and Block Locality}
\index{composite band!frame and alignment|textbf}
\index{gauge alignment!composite band|textbf}
\index{block hopping!gauge dependence|textbf}

\label{ch:training-composite-band-gauge-alignment}

\chapterstatus{pedagogical derivation and synthetic computational example. This
chapter explains the mathematical chain exercised by a controlled
one-dimensional rank-two alignment calculation, but it reports no M2 numerical
outcome, material calculation, localization result, topological classification,
scientific validation, or uncertainty quantification.}

\section{Learning goals}

After working through this chapter, a reader should be able to

\begin{enumerate}
  \item distinguish a matrix-valued periodic parent from a continuum parent and
        from a reduced represented operator;
  \item define a composite retained space by its frame and projector while
        allowing internal band crossings;
  \item separate an external spectral gap from neighboring-frame overlap
        conditioning;
  \item derive the covariance of a projected Hamiltonian under a
        momentum-dependent unitary gauge;
  \item connect a Bloch frame, a Wannier basis, and matrix-valued hopping
        blocks without confusing that connection with localization;
  \item construct neighborwise polar transport and interpret reciprocal sewing
        and closure holonomy;
  \item derive pointwise and globally constrained unitary Procrustes alignment;
  \item explain why complete spectra can be gauge invariant while finite-range
        locality is gauge dependent; and
  \item interpret the Frobenius, singular-value, gap, Hermiticity, and spectral
        diagnostics used in a bounded multiband calculation.
\end{enumerate}

Chapter~\ref{ch:training-projection-retained-spaces} introduced frames,
projectors, and exact retained operators. Chapter~\ref{ch:training-isolated-band-fourier-reduction}
then considered a rank-one retained Hamiltonian, for which the represented
operator at each reciprocal coordinate is the scalar energy $E_n(k)$. A
composite retained space has rank greater than one. Its Hamiltonian is a matrix,
and that matrix cannot be discussed independently of the frame in which it is
written.

The new sequence is
\begin{equation}
  \begin{aligned}
    \text{periodic parent}
    &\longrightarrow \text{composite space}
     \longrightarrow \text{frame path},\\
    \text{frame path}
    &\longrightarrow \text{transport and alignment}
     \longrightarrow H_F(\kappa),\\
    H_F(\kappa)
    &\longrightarrow \text{block hoppings}
     \longrightarrow \text{finite range}.
  \end{aligned}
  \label{eq:training-composite-sequence}
\end{equation}
Each arrow carries information that a table of eigenvalues alone does not
contain.

\section{A finite matrix-valued periodic parent}
\index{block Hamiltonian!periodic parent}

Let $G>0$ be one reciprocal period and let
$\kappa=k/G\in[-1/2,1/2)$. Consider the finite Bloch matrix
\begin{equation}
  H_{\mathrm p}(\kappa)
  =\sum_{R\in\mathcal R_{\mathrm p}}
    e^{2\pi iR\kappa}T_R,
  \label{eq:training-composite-parent}
\end{equation}
where each $T_R\in\mathbb C^{D\times D}$ acts on a declared $D$-dimensional
orbital coordinate space. The parent is Hermitian at every $\kappa$ when
\begin{equation}
  T_{-R}=T_R^{\dagger}
  \label{eq:training-composite-parent-hermiticity}
\end{equation}
for every represented displacement, including $T_0=T_0^{\dagger}$. Indeed,
\begin{align}
  H_{\mathrm p}(\kappa)^{\dagger}
  &=\sum_R e^{-2\pi iR\kappa}T_R^{\dagger}
    =\sum_R e^{-2\pi iR\kappa}T_{-R}
    =H_{\mathrm p}(\kappa).
  \label{eq:training-composite-fiber-hermiticity}
\end{align}

This parent differs from the continuum cosine Hamiltonian used in
Chapter~\ref{ch:training-isolated-band-fourier-reduction}. Equation~\eqref{eq:training-composite-parent}
is already an exact finite represented model. Evaluating its finite sum does
not introduce a plane-wave cutoff or finite-difference stencil error. It may be
a useful controlled parent, but it is not thereby a continuum Hamiltonian or a
material model.

The symbols $R$ and $T_R$ already resemble lattice hoppings. They describe the
\emph{ambient parent} coordinate space. The reduced blocks constructed later
act in a different, rank-$M$ retained space. Equal notation or equal range does
not identify those two spaces.

\section{Composite retention and the external gap}
\index{composite band!external gap}
\index{spectral gap!external}

At each reciprocal coordinate, diagonalize the parent:
\begin{equation}
  H_{\mathrm p}(\kappa)u_n(\kappa)
  =E_n(\kappa)u_n(\kappa),
  \qquad
  E_1(\kappa)\leq\cdots\leq E_D(\kappa).
  \label{eq:training-composite-eigensystem}
\end{equation}
Suppose the lowest $M<D$ states are retained. Collect an orthonormal basis for
their span in the frame matrix
\begin{equation}
  F(\kappa)
  =\begin{bmatrix}u_1(\kappa)&\cdots&u_M(\kappa)\end{bmatrix},
  \qquad
  F(\kappa)^{\dagger}F(\kappa)=I_M.
  \label{eq:training-composite-frame}
\end{equation}
The corresponding ambient-space projector is
\begin{equation}
  P(\kappa)=F(\kappa)F(\kappa)^{\dagger}.
  \label{eq:training-composite-projector}
\end{equation}

For this lowest-$M$ selection, the direct external gap is
\begin{equation}
  \gamma_{\mathrm{ext}}(\kappa)
  =E_{M+1}(\kappa)-E_M(\kappa).
  \label{eq:training-composite-external-gap}
\end{equation}
A sampled calculation may retain
\begin{equation}
  \gamma_{\mathrm{ext}}^{\mathrm{sampled}}
  =\min_{\kappa\in\mathcal K_{\mathrm{declared}}}
   \gamma_{\mathrm{ext}}(\kappa).
  \label{eq:training-composite-sampled-gap}
\end{equation}
A positive sampled value supports separation on the declared coordinates. It
does not prove a positive gap between those coordinates.

The retained bands may cross one another internally. Such an internal
degeneracy can make individual eigenvectors or energy-ordered labels unstable
without destroying the $M$-dimensional projector. This is a principal reason to
retain a composite space: the projector can remain smooth and meaningful even
when the individual columns returned by an eigensolver are not.

An energy gap and a neighboring-frame overlap diagnose different things. The
gap separates retained and excluded eigenvalues at one coordinate. An overlap
singular value, introduced below, measures how well two sampled retained spaces
see one another. Neither quantity can replace the other.

\section{Frames, gauges, and projected Hamiltonians}
\index{gauge!composite frame}
\index{projected Hamiltonian!gauge covariance}

For any $U(\kappa)\in U(M)$, define another frame
\begin{equation}
  F'(\kappa)=F(\kappa)U(\kappa).
  \label{eq:training-composite-gauge-change}
\end{equation}
Its projector is unchanged:
\begin{equation}
  F'F'^{\dagger}=FUU^{\dagger}F^{\dagger}=FF^{\dagger}=P.
  \label{eq:training-composite-projector-invariance}
\end{equation}
The multiplication order matters: the frame is $D\times M$, whereas $U$ is
$M\times M$.

The parent compressed into the frame is
\begin{equation}
  H_F(\kappa)
  =F(\kappa)^{\dagger}H_{\mathrm p}(\kappa)F(\kappa)
  \in\mathbb C^{M\times M}.
  \label{eq:training-composite-projected-operator}
\end{equation}
Under Equation~\eqref{eq:training-composite-gauge-change},
\begin{equation}
  H_{F'}(\kappa)
  =U(\kappa)^{\dagger}H_F(\kappa)U(\kappa).
  \label{eq:training-composite-operator-covariance}
\end{equation}
Thus the two matrices generally differ entry by entry even though they represent
the same retained abstract operator. Their eigenvalues agree. Subtracting them
before alignment measures a coordinate change, not a physical perturbation.

If the columns of $F$ are the individual eigenvectors, then $H_F$ is diagonal.
A smooth composite gauge may mix those eigenvectors, in which case $H_F$ need
not be diagonal. Diagonality and smoothness are different coordinate
properties.

\section{The Bloch--Wannier bridge}
\index{Wannier transformation!composite frame}
\index{Bloch--Wannier bridge}

Let $\{\lvert\psi_{m k}\rangle\}_{m=1}^{M}$ be a retained Bloch frame on a
uniform $N_k$-point mesh. For a chosen unitary gauge $U(k)$, define
\begin{equation}
  \lvert w_{nR}\rangle
  =\frac{1}{\sqrt{N_k}}
   \sum_k e^{-ikRa}
   \sum_{m=1}^{M}
   \lvert\psi_{mk}\rangle U_{mn}(k),
  \label{eq:training-composite-wannier-transform}
\end{equation}
where $R$ is an integer cell displacement. This is the finite
Bloch--Wannier transformation~\cite{wannier1937,marzari1997,marzari2012}.
It constructs a basis only after a frame and its boundary convention have been
chosen.

Let $H_U(k)=U(k)^{\dagger}H_F(k)U(k)$. Orthogonality between different Bloch
coordinates gives
\begin{align}
  \langle w_{m0}\rvert\widehat H\lvert w_{nR}\rangle
  &=\frac{1}{N_k}\sum_k
    e^{-ikRa}[H_U(k)]_{mn}.
  \label{eq:training-composite-wannier-hopping}
\end{align}
The inverse finite transform is
\begin{equation}
  H_U(k)
  =\sum_R e^{ikRa}H_U(R),
  \qquad
  [H_U(R)]_{mn}
  =\langle w_{m0}\rvert\widehat H\lvert w_{nR}\rangle.
  \label{eq:training-composite-wannier-inverse}
\end{equation}
Equations~\eqref{eq:training-composite-wannier-hopping} and
\eqref{eq:training-composite-wannier-inverse} are the missing bridge between a
Bloch-frame Hamiltonian and block hoppings.

For one isolated eigenband, $M=1$ and $H_F(k)=E_n(k)$. The phase $U(k)$ cancels
from $U(k)^*E_n(k)U(k)$, so the complete scalar hopping coefficients depend only
on the dispersion. The Wannier function itself still depends on the phase. For
$M>1$, a momentum-dependent $U(k)$ generally changes the matrix
$H_U(k)$ and therefore its individual real-space blocks.

The transform of Hamiltonian matrices is not a localization calculation. It
does not by itself calculate Wannier functions, centers, spreads, probability
densities, or a minimizer of a localization functional. Those quantities
require eigenvectors, overlaps, and an explicit localization construction.

\section{Comparing retained frames}
\index{frame defect}
\index{projector defect}
\index{principal angle!retained spaces}

Let $F_{\mathrm r}$ and $F_{\mathrm c}$ be reference and candidate frames in
the same ambient coordinates. Three comparisons must remain distinct.

The raw frame defect
\begin{equation}
  d_{\mathrm{frame}}
  =\lVert F_{\mathrm c}-F_{\mathrm r}\rVert_{\mathrm F}
  \label{eq:training-composite-frame-defect}
\end{equation}
is gauge dependent. The projector defect
\begin{equation}
  d_{\mathrm{proj}}
  =\lVert
     F_{\mathrm c}F_{\mathrm c}^{\dagger}
     -F_{\mathrm r}F_{\mathrm r}^{\dagger}
   \rVert_{\mathrm F}
  \label{eq:training-composite-projector-defect}
\end{equation}
is invariant under independent unitary rotations of either frame. It vanishes
exactly when the two orthogonal projectors agree.

Finally, the singular values of
\begin{equation}
  S=F_{\mathrm c}^{\dagger}F_{\mathrm r}
  \label{eq:training-composite-frame-overlap}
\end{equation}
are the cosines of the principal angles between the retained spaces. Values
near one indicate compatible subspaces; values near zero expose directions
that cannot be stably identified by a unitary rotation. Equal rank alone does
not guarantee favorable singular values.

The Frobenius norm used here is
\begin{equation}
  \lVert A\rVert_{\mathrm F}
  =\left(\sum_{i,j}|A_{ij}|^2\right)^{1/2}
  =\sqrt{\operatorname{tr}(A^{\dagger}A)}.
  \label{eq:training-composite-frobenius}
\end{equation}
A frame or projector defect is dimensionless. A Hamiltonian-matrix defect has
energy units. The same norm formula does not erase that dimensional
difference.

\section{Neighborwise polar transport}
\index{parallel transport!polar}
\index{singular value!neighboring frames}

Raw eigensolver frames at neighboring points can contain arbitrary phases and
unitary mixtures. Suppose a transported frame $F_j$ has already been chosen and
$G_{j+1}$ is the next raw frame. Form
\begin{equation}
  S_j=F_j^{\dagger}G_{j+1}.
  \label{eq:training-composite-neighbor-overlap}
\end{equation}
Let its singular-value decomposition be
\begin{equation}
  S_j=L_j\Sigma_jR_j^{\dagger}.
  \label{eq:training-composite-neighbor-svd}
\end{equation}
Then choose
\begin{equation}
  Q_j=R_jL_j^{\dagger},
  \qquad
  F_{j+1}=G_{j+1}Q_j.
  \label{eq:training-composite-polar-transport}
\end{equation}
The new overlap is
\begin{equation}
  F_j^{\dagger}F_{j+1}
  =L_j\Sigma_jL_j^{\dagger},
  \label{eq:training-composite-positive-overlap}
\end{equation}
which is Hermitian positive semidefinite. This is the unitary polar or
orthogonal-Procrustes step that makes the next frame as close as possible to the
previous one in the declared Frobenius metric~\cite{schonemann1966,higham1986}.

The minimum singular value over neighbor and boundary overlaps is a conditioning
diagnostic. A frozen threshold can make insufficient overlap a fail-closed stop
condition. It is not a universal physical constant, and passing it does not
prove continuous smoothness between mesh points.

\subsection{Sewing and closure holonomy}
\index{reciprocal sewing}
\index{holonomy!closure}

The first and last stored mesh points are not duplicate endpoints. To close the
reciprocal loop, the first frame must be mapped across the zone boundary by a
declared ambient sewing map $S_G$. After neighborwise transport, form
\begin{equation}
  C=F_{N_k-1}^{\dagger}S_GF_0.
  \label{eq:training-composite-closure-overlap}
\end{equation}
The unitary polar factor of $C$ is the discrete closure holonomy. If its
eigenvalues are $e^{i\phi_1},\ldots,e^{i\phi_M}$, the phases are defined only
modulo $2\pi$ and their ordering can change at degeneracies or branch cuts.
They must therefore be treated as a phase set with an explicit branch
convention~\cite{marzari2012}.

A periodic orbital-coordinate parent may use $S_G=I$ because its ambient basis
is exactly periodic. Other representations can require a nontrivial sewing map.
For example, reciprocal-index plane-wave coordinates relabel across
$k\mapsto k+G$. Identity sewing must never be assumed merely because the
Hamiltonian spectrum is periodic.

One can distribute a matrix root of the closure holonomy along the path to
produce a closed smooth representative. This is a gauge construction. Its
closure eigenphases are Wilson-loop-like data, but a single finite loop does not
by itself establish a polarization, a topological invariant, or an
infinite-mesh localization theorem.

\section{A controlled momentum-dependent gauge attack}
\index{gauge attack!rank two}

For a rank-two retained space, a transparent periodic rotation is
\begin{equation}
  A(\kappa)
  =e^{-i\theta(\kappa)\sigma_y}
  =\begin{pmatrix}
      \cos\theta(\kappa)&-\sin\theta(\kappa)\\
      \sin\theta(\kappa)& \cos\theta(\kappa)
    \end{pmatrix},
  \label{eq:training-composite-attack}
\end{equation}
with
\begin{equation}
  \theta(\kappa)
  =\theta_0+\sum_{m=1}^{M_a}a_m\sin(2\pi m\kappa).
  \label{eq:training-composite-attack-angle}
\end{equation}
Integer harmonics make the attack periodic. Define
\begin{equation}
  F_{\mathrm a}(\kappa)=F_{\mathrm r}(\kappa)A(\kappa).
  \label{eq:training-composite-attacked-frame}
\end{equation}
Then
\begin{equation}
  P_{\mathrm a}=P_{\mathrm r},
  \qquad
  H_{\mathrm a}=A^{\dagger}H_{\mathrm r}A.
  \label{eq:training-composite-attack-invariants}
\end{equation}
The frame and matrix entries can change substantially while the projector and
pointwise spectrum remain unchanged.

A finite Fourier series for $\theta$ does not imply that the entries of
$A(\kappa)$ have the same finite Fourier range: the nonlinear sine and cosine
of $\theta(\kappa)$ can generate longer Fourier tails. Consequently a smooth,
known gauge transformation can change the apparent range of the projected
Hamiltonian without changing its complete spectral content.

This attack is a controlled positive test. Recovering it does not show that an
unknown alignment between independent physical calculations has been found.

\section{Pointwise and globally constrained alignment}
\index{Procrustes alignment!pointwise}
\index{alignment!global unitary}

\subsection{Pointwise unitary Procrustes alignment}

At each reciprocal coordinate, solve
\begin{equation}
  \min_{U(\kappa)\in U(M)}
  \lVert F_{\mathrm c}(\kappa)U(\kappa)
          -F_{\mathrm r}(\kappa)\rVert_{\mathrm F}.
  \label{eq:training-composite-pointwise-procrustes}
\end{equation}
If
\begin{equation}
  F_{\mathrm c}^{\dagger}F_{\mathrm r}
  =L\Sigma R^{\dagger},
\end{equation}
then one minimizing rotation is
\begin{equation}
  U_{\star}=LR^{\dagger}.
  \label{eq:training-composite-pointwise-solution}
\end{equation}
For the constructed attack $F_{\mathrm c}=F_{\mathrm r}A$, the overlap is
$A^{\dagger}$ and pointwise alignment recovers $A^{\dagger}$ up to numerical
and nonuniqueness effects. Exact recovery in that setting is expected
mathematics, not evidence for a general optimizer.

Pointwise freedom is a large admissible family: every reciprocal coordinate
receives its own unitary. It can be inappropriate when the intended physical
comparison permits only a smaller parameterized transformation.

\subsection{One global unitary}

A deliberately restricted comparison uses one $U\in U(M)$ for the entire path:
\begin{equation}
  \min_{U\in U(M)}
  \sum_j\lVert F_{\mathrm c}(\kappa_j)U
                   -F_{\mathrm r}(\kappa_j)\rVert_{\mathrm F}^2.
  \label{eq:training-composite-global-procrustes}
\end{equation}
Let
\begin{equation}
  B=\sum_jF_{\mathrm c}(\kappa_j)^{\dagger}
          F_{\mathrm r}(\kappa_j)
   =L_B\Sigma_BR_B^{\dagger}.
  \label{eq:training-composite-global-aggregate}
\end{equation}
A minimizing global rotation is
\begin{equation}
  U_{\mathrm{global}}=L_BR_B^{\dagger}.
  \label{eq:training-composite-global-solution}
\end{equation}
If $B$ is rank deficient or has repeated singular structure, the minimizer need
not be unique or stable. Singular values and residuals should therefore be
retained when uniqueness matters.

One constant unitary generally cannot undo a genuinely momentum-dependent
attack. A nonzero global residual establishes a limitation of this restricted
family for the declared data. It does not establish failure of every
nonidentity alignment family. Enlarging the admissible family is a new model
class and requires a separately frozen calculation.

This distinction later becomes part of an adequacy decision. A model parameter
and an alignment variable play different roles: the first selects a candidate
Hamiltonian, while the second identifies its coordinates with the target
within a declared family. Chapter~\ref{ch:admissible-sets-certificates} defines
operator-admissible parameter sets by minimizing only over that frozen
alignment family. Failure for one global unitary, a finite list of rotations,
or a pointwise family therefore produces three different claims.

\section{Complete block hoppings and gauge-dependent locality}
\index{block hopping!complete transform}
\index{locality!gauge dependence}

On the complete centered mesh, transform the projected matrices according to
\begin{equation}
  H_F(R)
  =\frac{1}{N_k}\sum_{j=0}^{N_k-1}
   e^{-2\pi iR\kappa_j}H_F(\kappa_j).
  \label{eq:training-composite-block-transform}
\end{equation}
The inverse is
\begin{equation}
  H_F(\kappa_j)
  =\sum_{R\in\mathcal R_{N_k}}
   e^{2\pi iR\kappa_j}H_F(R).
  \label{eq:training-composite-block-inverse}
\end{equation}
Hermiticity requires the modular relation
\begin{equation}
  H_F(\overline{-R})=H_F(R)^{\dagger},
  \label{eq:training-composite-block-hermiticity}
\end{equation}
including the self-opposite even-mesh Nyquist representative.

A symmetric range-$R_{\max}$ model retains $|R|\leq R_{\max}$. Its omitted-block
norm is
\begin{equation}
  \eta_F(R_{\max})
  =\left(
    \sum_{|R|>R_{\max}}\lVert H_F(R)\rVert_{\mathrm F}^2
   \right)^{1/2}.
  \label{eq:training-composite-omitted-block-norm}
\end{equation}
This quantity has energy units. It measures represented block content outside
the selected range; it is not a probability or uncertainty.

For a complete transform, reference, attacked, and correctly realigned gauges
have the same pointwise eigenvalues. After truncation they need not. The
momentum-dependent gauge redistributes matrix content among the $R$ blocks, so
it changes both $\eta_F$ and the spectral error of a finite-range model. This is
why localization or gauge selection matters to short-range approximation even
though it does not change the complete retained spectrum.

The maximum spectral error on a declared set $\mathcal K$ is
\begin{equation}
  \epsilon_{\mathrm{spec},\max}
  =\max_{\kappa\in\mathcal K}
    \max_{1\leq n\leq M}
    \left|
      E_n^{\mathrm{truncated}}(\kappa)-E_n^{\mathrm{target}}(\kappa)
    \right|.
  \label{eq:training-composite-spectral-error}
\end{equation}
It is an eigenvalue metric. It does not measure entrywise or aligned-operator
agreement.

\section{Training and withheld roles}
\index{withheld mesh!multiband alignment}

The complete training mesh determines raw eigenspaces, transported frames,
alignment maps, projected matrices, Fourier blocks, and every finite-range
candidate. A staggered disjoint mesh can then evaluate spectra of those frozen
candidates between training coordinates.

In a bounded composite-space protocol, withheld coordinates may also enter a
prospectively declared sampled external-gap gate. That use must be named: the
withheld values then contribute to a frozen prerequisite and to final spectral
evaluation, even though they do not alter frames, coefficients, ranges, or
tolerances. They must not be used after inspection to redesign the attack or
select which range is reported while retaining an untouched-withheld label.

No withheld frame is needed merely to compare the eigenvalues of an
interpolated block Hamiltonian with the parent spectrum. Consequently such a
comparison does not establish withheld frame, projector, or operator-matrix
agreement.

\section{A diagnostic ledger}

\begin{table}[htbp]
  \centering
  \caption{Distinct diagnostics in a bounded composite-frame calculation.}
  \label{tab:training-composite-diagnostics}
  \small
  \begin{tabular}{@{}p{0.25\textwidth}p{0.29\textwidth}p{0.13\textwidth}p{0.23\textwidth}@{}}
    \toprule
    Quantity & Compared objects & Unit & Interpretation \\
    \midrule
    External gap & retained and excluded eigenvalues & energy & sampled spectral separation \\
    Minimum overlap singular value & neighboring or closure frames & unitless & transport conditioning \\
    Closure eigenphases & unitary closure holonomy & radians & branch-dependent phase set \\
    Frame Frobenius defect & two coordinate frames & unitless & gauge-dependent basis difference \\
    Projector Frobenius defect & two retained subspaces & unitless & gauge-invariant subspace difference \\
    Projected-operator defect & two matrices in declared coordinates & energy & coordinate-sensitive operator difference \\
    Block Hermiticity defect & $H(-R)$ and $H(R)^{\dagger}$ & energy & represented structural consistency \\
    Omitted-block norm & discarded hopping blocks & energy & range-truncation diagnostic \\
    Maximum spectral error & ordered eigenvalues on a mesh & energy & worst sampled band discrepancy \\
    \bottomrule
  \end{tabular}
\end{table}

Each tolerance belongs to one row and one predicate. A dimensionless frame
tolerance cannot silently become an energy tolerance. A numerical consistency
tolerance is not a physical uncertainty or scientific acceptance threshold.

\section{A transparent computational example}

The following synthetic example constructs a rank-two retained frame, applies a
known periodic gauge attack, recovers it pointwise, and compares complete and
truncated block representations. It is deliberately small and contains no
material parameters.

\begin{pythonexample}{aligning a synthetic composite Bloch frame}
import numpy as np
from scipy.linalg import schur

# Define a Hermitian four-orbital parent with cell ranges -1, 0, and +1.
onsite = np.array(
    [
        [-1.20, 0.06, 0.00, 0.00],
        [0.06, -0.35, 0.00, 0.00],
        [0.00, 0.00, 1.00, 0.04],
        [0.00, 0.00, 0.04, 1.75],
    ],
    dtype=np.complex128,
)
positive = np.array(
    [
        [0.15, 0.00, 0.12, 0.00],
        [0.00, -0.10, 0.00, 0.10],
        [0.00, 0.00, 0.05, 0.00],
        [0.00, 0.00, 0.00, -0.04],
    ],
    dtype=np.complex128,
)
parent_R = np.array([-1, 0, 1], dtype=np.int64)
parent_blocks = np.array(
    [positive.conjugate().T, onsite, positive],
    dtype=np.complex128,
)

# Evaluate the exact finite parent on a complete centered reciprocal mesh.
N = 16
kappa = -0.5 + np.arange(N, dtype=np.float64) / N
parent_phases = np.exp(2j * np.pi * np.outer(kappa, parent_R))
parent = np.einsum("kr,rij->kij", parent_phases, parent_blocks)

# Retain the lowest two eigenvectors at every reciprocal coordinate.
raw_frames = []
retained_energies = []
for matrix in parent:
    values, vectors = np.linalg.eigh(matrix)
    retained_energies.append(values[:2])
    raw_frames.append(vectors[:, :2])
raw_frames = np.asarray(raw_frames, dtype=np.complex128)
retained_energies = np.asarray(retained_energies, dtype=np.float64)

# Apply neighborwise polar transport to remove arbitrary eigensolver gauges.
transported = [raw_frames[0].copy()]
neighbor_singular_values = []
for index in range(N - 1):
    overlap = transported[index].conjugate().T @ raw_frames[index + 1]
    left, singular_values, right_h = np.linalg.svd(overlap)
    neighbor_singular_values.extend(singular_values)
    rotation = right_h.conjugate().T @ left.conjugate().T
    transported.append(raw_frames[index + 1] @ rotation)

# Close the half-open reciprocal loop with identity ambient sewing.
closure = transported[-1].conjugate().T @ transported[0]
left, closure_singular_values, right_h = np.linalg.svd(closure)
neighbor_singular_values.extend(closure_singular_values)
unitary_closure = left @ right_h
triangular, eigenvectors = schur(unitary_closure, output="complex")
closure_phases = np.angle(np.diag(triangular))
order = np.argsort(closure_phases)
closure_phases = closure_phases[order]
eigenvectors = eigenvectors[:, order]

# Distribute the closure holonomy over the represented loop.
for index in range(N):
    fraction = index / N
    root = (
        eigenvectors
        @ np.diag(np.exp(1j * closure_phases * fraction))
        @ eigenvectors.conjugate().T
    )
    transported[index] = transported[index] @ root
reference = np.asarray(transported, dtype=np.complex128)

# Apply a known periodic, momentum-dependent rank-two gauge rotation.
theta = 0.30 + 0.65 * np.sin(2.0 * np.pi * kappa)
attacks = np.array(
    [
        [[np.cos(value), -np.sin(value)],
         [np.sin(value), np.cos(value)]]
        for value in theta
    ],
    dtype=np.complex128,
)
attacked = np.einsum("kdi,kij->kdj", reference, attacks)

# Projectors are invariant even though the frame matrices differ.
reference_projectors = np.einsum(
    "kdi,kei->kde", reference, reference.conjugate()
)
attacked_projectors = np.einsum(
    "kdi,kei->kde", attacked, attacked.conjugate()
)
np.testing.assert_allclose(attacked_projectors, reference_projectors, atol=1e-13)

# Recover one unitary Procrustes rotation independently at each mesh point.
aligned = []
recovered_rotations = []
for reference_frame, candidate_frame in zip(reference, attacked, strict=True):
    overlap = candidate_frame.conjugate().T @ reference_frame
    left, _, right_h = np.linalg.svd(overlap)
    rotation = left @ right_h
    recovered_rotations.append(rotation)
    aligned.append(candidate_frame @ rotation)
aligned = np.asarray(aligned, dtype=np.complex128)
recovered_rotations = np.asarray(recovered_rotations, dtype=np.complex128)
np.testing.assert_allclose(aligned, reference, atol=1e-13)
np.testing.assert_allclose(
    recovered_rotations,
    attacks.conjugate().transpose(0, 2, 1),
    atol=1e-13,
)

# Fit the distinct restricted family containing one unitary for the whole path.
aggregate = sum(
    candidate.conjugate().T @ target
    for candidate, target in zip(attacked, reference, strict=True)
)
left, _, right_h = np.linalg.svd(aggregate)
global_rotation = left @ right_h
globally_aligned = np.asarray(
    [frame @ global_rotation for frame in attacked],
    dtype=np.complex128,
)
global_frame_defect = max(
    np.linalg.norm(candidate - target)
    for candidate, target in zip(globally_aligned, reference, strict=True)
)

# Compress the same parent into the frame choices.
def project(frames):
    return np.asarray(
        [
            frame.conjugate().T @ matrix @ frame
            for matrix, frame in zip(parent, frames, strict=True)
        ],
        dtype=np.complex128,
    )

reference_operator = project(reference)
attacked_operator = project(attacked)
aligned_operator = project(aligned)
global_operator = project(globally_aligned)

# All complete projected operators have the retained parent eigenvalues.
for operator in (
    reference_operator,
    attacked_operator,
    aligned_operator,
    global_operator,
):
    values = np.linalg.eigvalsh(operator)
    np.testing.assert_allclose(values, retained_energies, atol=1e-13)

# Transform each matrix path into complete block hoppings.
cell_R = np.arange(-N // 2, N // 2, dtype=np.int64)
def blocks_from_samples(samples):
    return np.asarray(
        [
            np.mean(
                samples
                * np.exp(-2j * np.pi * R * kappa)[:, None, None],
                axis=0,
            )
            for R in cell_R
        ],
        dtype=np.complex128,
    )

reference_blocks = blocks_from_samples(reference_operator)
attacked_blocks = blocks_from_samples(attacked_operator)
aligned_blocks = blocks_from_samples(aligned_operator)
np.testing.assert_allclose(aligned_blocks, reference_blocks, atol=1e-13)

# A fixed finite range can retain different amounts in different gauges.
range_mask = np.abs(cell_R) <= 1
reference_omitted = np.linalg.norm(reference_blocks[~range_mask])
attacked_omitted = np.linalg.norm(attacked_blocks[~range_mask])

print("minimum neighbor/closure singular value:", min(neighbor_singular_values))
print("one-global-unitary frame defect:", global_frame_defect)
print("reference omitted-block norm:", reference_omitted)
print("attacked omitted-block norm:", attacked_omitted)
\end{pythonexample}

The example establishes finite-array identities for a synthetic parent. It does
not establish continuous-gap protection, optimal localization, or a physical
alignment. The default matrix norm used by \texttt{numpy.linalg.norm} is the
Frobenius norm in these calls; production interfaces should name that choice
rather than rely on an unstated default.

\documentcallout{Notebook to do}{Create
\path{examples/tutorials/research-monograph/foundations/composite_band_gauge_alignment.ipynb}
only after this lecture contract is stable. Include frame/projector invariance,
neighbor singular values, polar transport, reciprocal sewing, closure phases,
pointwise and global Procrustes alignment, block Hermiticity, gauge-dependent
truncation, and disjoint withheld spectral evaluation. Keep the example
synthetic and do not create or execute the notebook as part of the present
manuscript change.}

\section{Common mistakes}

\begin{enumerate}
  \item Tracking individual energy-ordered eigenvectors through an internal
        degeneracy instead of retaining the composite projector.
  \item Treating an external energy gap as a substitute for overlap singular
        values, or vice versa.
  \item Comparing frame matrices before accounting for unitary gauge freedom.
  \item Calling a large unaligned matrix defect a physical operator error.
  \item Assuming periodic eigenvalues imply identity reciprocal sewing in every
        basis.
  \item Comparing closure eigenphases as ordinary real numbers without branch
        and permutation handling.
  \item Describing exact recovery of a constructed pointwise attack as a
        general alignment solution.
  \item Concluding from failure of one global unitary that every nonidentity
        alignment family fails.
  \item Assuming a finite harmonic series for the attack angle gives a
        finite-range attacked Hamiltonian.
  \item Calling a smooth transported gauge maximally localized without
        evaluating a localization functional.
  \item Inferring Wannier centers, spreads, or densities from Hamiltonian
        blocks alone.
  \item Assuming gauge invariance of the complete spectrum survives an
        arbitrary finite-range truncation.
  \item Treating an omitted-block Frobenius norm as a probability or physical
        uncertainty.
  \item Using withheld spectra to redesign the gauge or range while preserving
        an untouched-evaluation claim.
\end{enumerate}

\section{Exercises}

\begin{enumerate}
  \item Starting from Equation~\eqref{eq:training-composite-parent}, prove
        Equation~\eqref{eq:training-composite-fiber-hermiticity}.
  \item Verify the dimensions of every factor in
        Equations~\eqref{eq:training-composite-projector} and
        \eqref{eq:training-composite-projected-operator}.
  \item Prove projector invariance and projected-operator covariance under
        $F\mapsto FU$.
  \item Derive Equation~\eqref{eq:training-composite-wannier-hopping} from the
        finite Wannier transformation and Bloch orthogonality.
  \item Show explicitly why the gauge cancels from the complete scalar
        isolated-band Hamiltonian but not from a general composite matrix.
  \item Derive the positive overlap in
        Equation~\eqref{eq:training-composite-positive-overlap}.
  \item Prove the pointwise Procrustes solution by maximizing
        $\operatorname{Re}\operatorname{tr}(U^{\dagger}F_{\mathrm c}^{\dagger}F_{\mathrm r})$.
  \item Derive the aggregate matrix in
        Equation~\eqref{eq:training-composite-global-aggregate}.
  \item Construct an attack that is constant in $\kappa$ and explain why the
        pointwise and global alignment families then contain the same inverse.
  \item Compare a frame defect, projector defect, and spectral defect for two
        frames related by a nontrivial unitary.
  \item Show that symmetric block truncation preserves
        Equation~\eqref{eq:training-composite-block-hermiticity} when the
        complete representative set is paired modulo $N_k$.
  \item Explain what additional evidence would be required to interpret closure
        eigenphases as Wannier centers or topological data.
\end{enumerate}

\section{Evidence boundary and next use}

The derivations in this chapter establish finite-dimensional identities under
stated frame, mesh, sewing, Fourier, and norm conventions. The Python example
checks those identities for synthetic arrays. Neither establishes that an
unknown physical alignment has been recovered, that a chosen gauge is
maximally localized, that a hopping range is adequate for an application, or
that a material is represented correctly.

A controlled M2-like calculation can use this chain to test a known gauge
attack, distinguish pointwise from globally constrained alignment, and expose
gauge-dependent finite-range locality. Such a calculation remains a bounded
positive and negative control. Its retained result and independent verifier own
any numerical outcomes; this chapter does not import them.

Part~II uses these foundations to ask whether an alignment or reduced model is
adequate for its declared purpose. Appendix~\ref{app:one-dimensional-reduction}
applies composite transport, closure, block hopping, and localization concepts
to retained controlled calculations. Later material chapters add projection
windows, disentanglement, geometry, orbital identity, spin, and energy-reference
requirements that the present synthetic construction deliberately omits.
